% GENERATED FILE. Edit canonical lessons, then run npm run generate:books.
% canonical-source-hash: fnv1a64:e6b0ad473580c8e7
% canonical-lessons: ZH-C494-gaokao, ZH-W494-kao-test, ZH-C494-kaosheng, ZH-C494-baokao, ZH-C494-kaoshang, ZH-C494-luobang

\chapter{The Entrance Exam}
\label{ch:zh-the-entrance-exam}
\hlchaptermodality{494}

\begin{chapteropening}
\textbf{By the end of this chapter:} \emph{I can talk about the college entrance exam, the candidates, applying, getting in and failing.}

Complete the last lesson of chapter 494: I can talk about the college entrance exam, the candidates, applying, getting in and failing.
\end{chapteropening}

\section[\texorpdfstring{\zh{高考} (gāokǎo)}{高考 (gāokǎo)}]{\zh{高考} (gāokǎo) — college entrance exam}

\label{lesson:ZH-C494-gaokao}



\begin{quote}
Before the new one: say the Mandarin for to tutor, then the Mandarin for to seek advice.
\end{quote}

\subsection*{You'll want to know: \zh{高考}}
\textbf{\zh{高考}} — \emph{gāokǎo} — \textquotedblleft{}college entrance exam\textquotedblright{}.

\begin{quote}
\emph{gāokǎo de shíhou wǒ gǎnjué hěn jǐnzhāng} — I felt very nervous during the entrance exam.
\end{quote}

The new character, in pieces: \zh{考} — six strokes: the top of \zh{老} \emph{lǎo}, four strokes, then a short bar and a stroke that drops, runs right, turns down and hooks. The next lesson puts it in your hand.

\begin{cousinweb}
\textbf{\zh{高考}} (\emph{gāokǎo}) is built from \zh{高} 'high' and \zh{考} 'to examine', both inherited from Middle Chinese, \emph{kaw} and \emph{khawX} (Guangyun, in Baxter's transcription), which glosses the second 'to check' and 'to complete'. It is short for a longer official name, the exam for the higher schools.
\end{cousinweb}

\begin{grammarlens}[title={What you've built}]
One more word for saying how it was.
\end{grammarlens}

\subsection*{Your turn}
\begin{itemize}
\raggedright
  \item \emph{Say it:} \emph{gāokǎo}
  \item \emph{Say it:} \emph{gāokǎo}, once more
  \item \emph{Say it:} \emph{gāokǎo de shíhou wǒ gǎnjué hěn jǐnzhāng} — I felt very nervous during the entrance exam.
  \item \emph{Recall:} say the Mandarin for to tutor, then the Mandarin for to seek advice, then say \emph{gāokǎo} again
\end{itemize}

\begin{tcolorbox}[breakable,colback=teal!4,colframe=teal!35!black,title={Before you move on}]
What does \zh{高考} mean? (\textquotedblleft{}College entrance exam\textquotedblright{}.) Say it once more.
\end{tcolorbox}
\section[\texorpdfstring{\zh{考} (kǎo)}{考 (kǎo)}]{\zh{考} — the top of lǎo, then two strokes}

\label{lesson:ZH-W494-kao-test}



\begin{quote}
The word first, then the character inside it.

\begin{itemize}
\raggedright
  \item \emph{Recall:} say \emph{gāokǎo} — \textbf{R1}, one lesson back, before it decays
\end{itemize}

\begin{quote}
\zh{老} \zh{高}
\end{quote}

\emph{Read it:} each one aloud, with nothing said first — \emph{lǎo}, \emph{gāo}
\end{quote}

\subsection*{Script — the top of lǎo, then two strokes}
\hlblockfigure{\detokenize{figures/ZH-W494-kao-test-filmstrip.pdf}}{How \zh{考} is written, stroke by stroke}

\begin{quote}
\zh{考}
\end{quote}

\textbf{The top of \zh{老}, four strokes}, as in \zh{老} and on the left of \zh{教}: a short bar, a vertical through it, a long bar, and a long stroke falling to the left.

\textbf{Then two strokes under it}: a short bar, and a stroke that starts under the bar's left end, drops a little, runs right, then turns down and hooks to the left.

4 + 2 = six.

\subsection*{Your turn}
\begin{itemize}
\raggedright
  \item \emph{Write it:} \zh{考}, all six, with the model in front of you
  \item \emph{Cover:} write it again from memory, and count the pieces, not the strokes
\end{itemize}

\begin{tcolorbox}[breakable,colback=teal!4,colframe=teal!35!black,title={Before you move on}]
How many strokes does \zh{考} have? (\textbf{Six}.) What are the first four? (\textbf{The top of \zh{老}}.) How does the last stroke end? (\textbf{It turns down and hooks to the left}.)

Source: \href{https://raw.githubusercontent.com/chanind/hanzi-writer-data/68d10a4b21150cae5e1ebbd223eed289cf32d90c/data/\%E8\%80\%83.json}{Hanzi Writer Data \zh{考}, pinned stroke paths and medians}.
\end{tcolorbox}
\section[\texorpdfstring{\zh{考生} (kǎoshēng)}{考生 (kǎoshēng)}]{\zh{考生} (kǎoshēng) — exam candidate}

\label{lesson:ZH-C494-kaosheng}



\begin{quote}
Before the new one: say the Mandarin for to seek advice, then the Mandarin for the college entrance exam.
\end{quote}

\subsection*{You'll want to know: \zh{考生}}
\textbf{\zh{考生}} — \emph{kǎoshēng} — \textquotedblleft{}exam candidate\textquotedblright{}.

\begin{quote}
\emph{kǎoshēng dōu hěn jǐnzhāng} — The candidates are all nervous.
\end{quote}

\begin{cousinweb}
\textbf{\zh{考生}} (\emph{kǎoshēng}) is built from \zh{考} 'to examine' and \zh{生} 'a student', both inherited from Middle Chinese, \emph{khawX} and \emph{srjaeng} (Guangyun, in Baxter's transcription), which glosses the second 'to grow'.
\end{cousinweb}

\begin{grammarlens}[title={What you've built}]
One more word for saying how it was.
\end{grammarlens}

\subsection*{Your turn}
\begin{itemize}
\raggedright
  \item \emph{Say it:} \emph{kǎoshēng}
  \item \emph{Say it:} \emph{kǎoshēng}, once more
  \item \emph{Say it:} \emph{kǎoshēng dōu hěn jǐnzhāng} — The candidates are all nervous.
  \item \emph{Recall:} say the Mandarin for to seek advice, then the Mandarin for the college entrance exam, then say \emph{kǎoshēng} again
\end{itemize}

\begin{tcolorbox}[breakable,colback=teal!4,colframe=teal!35!black,title={Before you move on}]
What does \zh{考生} mean? (\textquotedblleft{}Exam candidate\textquotedblright{}.) Say it once more.
\end{tcolorbox}
\section[\texorpdfstring{\zh{报考} (bàokǎo)}{报考 (bàokǎo)}]{\zh{报考} (bàokǎo) — to enter for an exam; to apply to a school by exam}

\label{lesson:ZH-C494-baokao}



\begin{quote}
Before the new one: say the Mandarin for the college entrance exam, then the Mandarin for an exam candidate.
\end{quote}

\subsection*{You'll want to know: \zh{报考}}
\textbf{\zh{报考}} — \emph{bàokǎo} — \textquotedblleft{}to enter for an exam; to apply to a school by exam\textquotedblright{}.

\begin{quote}
\emph{wǒ xiǎng bàokǎo Běijīng de dàxué} — I want to apply to universities in Beijing.
\end{quote}

\begin{cousinweb}
\textbf{\zh{报考}} (\emph{bàokǎo}) is built from \zh{报} 'to report' and \zh{考} 'to examine', both inherited from Middle Chinese, \emph{pawH} and \emph{khawX} (Guangyun, in Baxter's transcription), which glosses the first 'to report'.
\end{cousinweb}

\begin{grammarlens}[title={What you've built}]
One more word for saying how it was.
\end{grammarlens}

\subsection*{Your turn}
\begin{itemize}
\raggedright
  \item \emph{Say it:} \emph{bàokǎo}
  \item \emph{Say it:} \emph{bàokǎo}, once more
  \item \emph{Say it:} \emph{wǒ xiǎng bàokǎo Běijīng de dàxué} — I want to apply to universities in Beijing.
  \item \emph{Recall:} say the Mandarin for the college entrance exam, then the Mandarin for an exam candidate, then say \emph{bàokǎo} again
\end{itemize}

\begin{tcolorbox}[breakable,colback=teal!4,colframe=teal!35!black,title={Before you move on}]
What does \zh{报考} mean? (\textquotedblleft{}To enter for an exam; to apply to a school by exam\textquotedblright{}.) Say it once more.
\end{tcolorbox}
\section[\texorpdfstring{\zh{考上} (kǎoshàng)}{考上 (kǎoshàng)}]{\zh{考上} (kǎoshàng) — to get into (a school) by exam; to pass the entrance exam for}

\label{lesson:ZH-C494-kaoshang}



\begin{quote}
Before the new one: say the Mandarin for an exam candidate, then the Mandarin for to enter for an exam.
\end{quote}

\subsection*{You'll want to know: \zh{考上}}
\textbf{\zh{考上}} — \emph{kǎoshàng} — \textquotedblleft{}to get into (a school) by exam; to pass the entrance exam for\textquotedblright{}.

\begin{quote}
\emph{tā fùxí le yì nián cái kǎoshàng dàxué} — He revised for a year before he got into university.
\end{quote}

\begin{cousinweb}
\textbf{\zh{考上}} (\emph{kǎoshàng}) is built from \zh{考} 'to examine' and \zh{上} 'to go up', both inherited from Middle Chinese, \emph{khawX} and \emph{dzyangX} (Guangyun, in Baxter's transcription), which glosses the second 'to climb'.
\end{cousinweb}

\begin{grammarlens}[title={What you've built}]
One more word for saying how it was.
\end{grammarlens}

\subsection*{Your turn}
\begin{itemize}
\raggedright
  \item \emph{Say it:} \emph{kǎoshàng}
  \item \emph{Say it:} \emph{kǎoshàng}, once more
  \item \emph{Say it:} \emph{tā fùxí le yì nián cái kǎoshàng dàxué} — He revised for a year before he got into university.
  \item \emph{Recall:} say the Mandarin for an exam candidate, then the Mandarin for to enter for an exam, then say \emph{kǎoshàng} again
\end{itemize}

\begin{tcolorbox}[breakable,colback=teal!4,colframe=teal!35!black,title={Before you move on}]
What does \zh{考上} mean? (\textquotedblleft{}To get into (a school) by exam; to pass the entrance exam for\textquotedblright{}.) Say it once more.
\end{tcolorbox}
\section[luòbǎng]{luòbǎng — to fail an entrance exam}

\label{lesson:ZH-C494-luobang}



\begin{quote}
Before the new one: say the Mandarin for to enter for an exam, then the Mandarin for to pass the entrance exam for.
\end{quote}

\subsection*{You'll want to know: luòbǎng}
\emph{luòbǎng} — \textquotedblleft{}to fail an entrance exam\textquotedblright{}.

\begin{quote}
\emph{luòbǎng yǐhòu tā gǎnjué hěn nánguò} — After he failed the entrance exam, he felt awful.
\end{quote}

Its characters come later; for now it is a word for the ear.

\begin{cousinweb}
\emph{luòbǎng} is built from \emph{luò} 'to fall' and \emph{bǎng} 'a board', both inherited from Middle Chinese, \emph{lak} and \emph{pangX} (Guangyun, in Baxter's transcription), which glosses the first as leaves falling and the second 'a wooden board'. Falling off the list of names posted on a board is the usual explanation, but the dictionary does not give that sense.
\end{cousinweb}

\begin{grammarlens}[title={What you've built}]
One more word for saying how it was.
\end{grammarlens}

\subsection*{Your turn}
\begin{itemize}
\raggedright
  \item \emph{Say it:} \emph{luòbǎng}
  \item \emph{Say it:} \emph{luòbǎng}, once more
  \item \emph{Say it:} \emph{luòbǎng yǐhòu tā gǎnjué hěn nánguò} — After he failed the entrance exam, he felt awful.
  \item \emph{Recall:} say the Mandarin for to enter for an exam, then the Mandarin for to pass the entrance exam for, then say \emph{luòbǎng} again
\end{itemize}

\begin{tcolorbox}[breakable,colback=teal!4,colframe=teal!35!black,title={Before you move on}]
What does luòbǎng mean? (\textquotedblleft{}To fail an entrance exam\textquotedblright{}.) Say it once more.
\end{tcolorbox}
